\documentclass[11pt]{article}
\usepackage[a4paper,margin=25mm]{geometry}
\usepackage{amsmath,amssymb}
\usepackage[T1]{fontenc}
\usepackage{lmodern}
\usepackage{hyperref}
\setlength{\parindent}{0pt}
\setlength{\parskip}{7pt}
\title{An algebraic power-tower value at an interior algebraic base}
\author{SucRunBug}
\date{2 October 2026}
\begin{document}
\maketitle
\textbf{Conjecture 00000000217 -- FALSE.}

\textbf{Filed claim addressed.} At every algebraic point in the interior of $[e^{-e},e^{1/e}]$, the convergent power-tower value is a Mahler S-number.

\section*{Proof}
Take the algebraic base $x=1$. Since $e>0$, $e^{-e}<1<e^{1/e}$, so it is strictly inside the convergence interval. Every finite tower at base one equals one, and its limit is
\[h(1)=1.\]
The same value is forced by the stated functional equation $h(1)=1^{h(1)}=1$.

The value one is algebraic over $\mathbb Q$: it is a root of $X-1$. In Mahler's classification algebraic numbers are A-numbers, while S-numbers are transcendental. Therefore $h(1)$ is not an S-number. This refutes the interior-value conjunct and hence the complete filed conjecture. It makes no assertion about the proposed behavior near the lower endpoint.

\section*{Formalization scope}
Lean proves the strict interval inequalities, convergence of the actual finite tower sequence to one, algebraicity of the base and non-transcendence of the value. The written argument uses only the defining requirement that an S-number be transcendental.

\textbf{Reproduce.} In \texttt{lean4/}, run \texttt{lake exe cache get},
\texttt{lake build}, then \texttt{lake env lean Check.lean}. Versions:
Lean/Mathlib \texttt{v4.33.0}. Independent finite checks are in
\texttt{reproduce.py}; they are not a substitute for the complete proof.

\textbf{Source.} \url{https://github.com/The-Last-Math-Competition/The-Last-Math-Competition/blob/28a22efac52c4a7abd4bba8a143b1c1a06b7e177/conjectures/00000000217.md}
\end{document}
